\documentclass{../Resume}

\newcommand{\Titulo}{Resumen de asignatura: Biotecnología en alimentos}
\newcommand{\Fecha}{2026}

\begin{document}
\Portada
\Indice
\fancystyle
\onehalfspacing

% --------------------------------------------------------------------------------------------------- %

\part{Fundamentos de biotecnología alimentaria y tendencias en el consumo alimentario}

    \chapter{Ciencias y tecnología de los alimentos}
    \chapter{Composición de los alimentos aditivos y adulterantes alimentarios}
    \chapter{Clasificación de los alimentos}
    \chapter{Innovación en alimentos}
    \chapter{Patentes de biotecnología alimentaria}

\part{Tecnologías en conservación de alimentos y aplicaciones biotecnológicas en la industria alimentaria}

    \chapter{Metabolitos microbianos como mecanismo básico en la transformación, fermentación y conservación de alimentos}
    \chapter{Producción microbiana de aditivos alimentarios}
    \chapter{Diseño de procesos y nuevos productos aplicados a la biotecnología alimentaria y las buenas prácticas de manufactura}
    \chapter{HPP Pasteurización en frio}
    \chapter{Enzimas en la industria alimentaria}
    \chapter{Sustancias bioactivas en alimentos}
    \chapter{Aplicación de las sustancias bioactivas en la industria alimentaria}

\part{Alimentos funcionales}

    \chapter{Introducción a los alimentos funcionales}
    \chapter{Microflora intestinal: Probióticos y prebióticos}
    \chapter{Alimentos funcionales}
    \chapter{Matrices orgánicas transgénicos con aplicación en la industria aliementaria}

\part{Herramientas ómicas}

    \chapter{Introducción a la genómica y sus aplicaciones en la industria alimentaria}
    \chapter{Introducción a la proteómica y sus aplicaciones en la industria alimentaria}
    \chapter{Introducción a la nitrigenómica y sus aplicaciones en la industria alimentaria}
    \chapter{Organismos genéticamente modificados: mejora genética de bacterias y levaduras utilizadas en fermentaciones}

% --- BIBLIOGRAFÍA --- %

\newpage

\bibliography{bib}

\end{document}